\printedpage{54}
% head: A. Grothendieck : Dérivateurs, Ch. XIII, Édition des Universités de Montpellier II et Paris VII
% foot: Version 25 janvier 2004

\entry{$L_\pi$}{77}
\entry{$\tilde{I}$}{79}
\entry{$\underline{M}$}{86}
\entry{$I_M$}{87}
\entry{$M^\pi$}{87}

\clearpage

\printedpage{55}
% head: A. Grothendieck : Dérivateurs, Ch. XIII, Édition des Universités de Montpellier II et Paris VII
% foot: Version 25 janvier 2004

\section*{Terminological index}

\entry{accessible (left closed subset)}{75}
\entry{admissible (on the left)}{59}
\entry{category of models}{7}
\entry{closed Quillen category}{25}
\entry{left closed, right closed}{2}
\entry{Quillen couple}{5}
\entry{weakly saturated}{7}
\entry{cardinal filtration}{77}
\entry{strict cardinal filtration}{86}
\entry{arrow couniversally in $W$}{9}
\entry{arrow universally in $W$}{10}
\entry{large relative to $\pi$}{77}
\entry{path object}{32}
\entry{cylinder object}{32}
\entry{$\pi$-accessible (subset of a category equipped with a strict cardinal filtration)}{86}
\entry{left, right pre-Q-category of models}{12}
\entry{left, right pre-Quillen}{14}
\entry{closed Quillen pre-quintuple}{21}
\entry{proper}{8}
\entry{Quillen factorization property}{5}
\entry{Q-category of models}{8}
\entry{right Q-category of models}{21}
\entry{``closed'' Q-category of models}{16}
\entry{left, right Q-homotopy}{30}
\entry{Quillenizable}{50}
\entry{Quillenizable on Diag}{51}
\entry{Quillenizable on Diag uniformly and left exactly}{57}
\entry{uniformly left Quillenizable on Diag}{57}
\entry{uniform right Quillenization on Diag}{58}
\entry{uniform left Quillenization on Diag}{57}
\entry{uniform and left exact Quillenization on Diag}{57}
\entry{well-adapted representation}{84}
\entry{perfectly adapted representation}{86}
\entry{strictly admissible}{61}
\entry{Quillen structure}{28}

\clearpage

\printedpage{56}
% head: A. Grothendieck : Dérivateurs, Ch. XIII, Édition des Universités de Montpellier II et Paris VII
% foot: Version 25 janvier 2004

\entry{exact Quillen structure}{27}
\entry{left exact, right exact Quillen structure}{27}
\entry{weak Quillen structure}{27}
\entry{biexact weak Quillen structure}{28}
\entry{closed, left closed, right closed weak Quillen structure}{28}
\entry{perfect Quillen structure}{28}
\entry{closed Quillen system}{26}
\entry{weak Quillen system}{26}
\entry{left exact, right exact weak Quillen system}{26}
\entry{well-adapted direct system}{84}
\entry{perfectly adapted direct system}{86}
\entry{closed Quillen triple}{25}
\entry{Quillen triple}{14}
\entry{$W$-cofibration}{8}
\entry{couniversal $W$-equivalence}{9}
\entry{$W$-fibration}{8}

\clearpage

\printedpage{57}
% head: A. Grothendieck : Dérivateurs, Ch. XIII, Édition des Universités de Montpellier II et Paris VII
% foot: Version 25 janvier 2004

\section*{List of axioms}

\noindent $\mathcal{M}$ category

\noindent $W \ \subset\ \Fl(\mathcal{M})$ , \qquad $C$, $F$, $TC$, $TF \ \subset\ \Fl(\mathcal{M})$

\begin{description}[leftmargin=4.5em,labelwidth=4em]
\item[M 0] The category $\mathcal{M}$ admits finite inverse limits and finite direct limits.

\item[M 1] The arrows belonging to $C$ have the left lifting property relative to those belonging to $TF$, and the arrows belonging to $F$ have the right lifting property relative to those belonging to $TC$.

\item[M 2] Every arrow of $\mathcal{M}$ decomposes into an arrow of $C$ followed by an arrow of $TF$, and into an arrow of $TC$ followed by an arrow of $F$.

\item[M 3] The subset $F$ of $\Fl(\mathcal{M})$ is stable under composition and base change, $C$ is stable under composition and cobase change, $F$ and $C$ contain the isomorphisms.

\item[M 4] The subset $TF$ of $\Fl(\mathcal{M})$ is stable under composition and base change, $TC$ is stable under composition and cobase change, $TF$ and $TC$ contain the isomorphisms.

\item[M 5] If in a commutative triangle of $\mathcal{M}$ two of the three arrows are in $W$, the same holds for the third ; $W$ contains the isomorphisms.

\item[M$'$] $TC \ \subset\ C \cap W$, $TF \ \subset\ F \cap W$.

\item[M$''$] The arrows belonging to $C$ are $W$-cofibrations, and the arrows belonging to $F$ are $W$-fibrations.

\item[M$'_{\mathbf{0}}$(g)] $TC = C \cap W$.

\item[M$'_{\mathbf{0}}$(d)] $TF = F \cap W$.

\item[M 2\,bis] $W = TF \circ TC$.
\end{description}

\clearpage

\printedpage{58}
% head: A. Grothendieck : Dérivateurs, Ch. XIII, Édition des Universités de Montpellier II et Paris VII
% foot: Version 25 janvier 2004

\setcounter{footnote}{0}
\renewcommand{\thefootnote}{\fnsymbol{footnote}}
\section*{Glossary}

\subsection*{Section 1}

\begin{itemize}
\item $\Phi \ \subset\ \Fl(\mathcal{M})$

$\Phi$ is \textit{left closed} : there exists \ $\Psi \ \subset\ \Fl(\mathcal{M})$ such that $\Phi$ consists of the arrows having the left lifting property relative to the arrows of $\Psi$.

\item $\Psi \ \subset\ \Fl(\mathcal{M})$

$\Psi$ is \textit{right closed} : there exists \ $\Phi \ \subset\ \Fl(\mathcal{M})$ such that $\Psi$ consists of the arrows having the right lifting property relative to the arrows of $\Phi$.

\item $\Phi$, $\Psi \ \subset\ \Fl(\mathcal{M})$

$(\Phi, \Psi)$ has the \textit{Quillen factorization property} :
\begin{itemize}
\item[a)] the arrows of $\Phi$ have the left lifting property relative to the arrows of $\Psi$ ;
\item[b)] every arrow $f$ of $\mathcal{M}$ factors as
\[
f = pi \quad , \qquad i \in \Phi \ , \quad p \in \Psi \ .
\]
\end{itemize}

\item $(\Phi, \Psi)$ \textit{Quillen couple} :
\begin{itemize}
\item[a)] $(\Phi, \Psi)$ has the Quillen factorization property ;
\item[b)] $\Phi$ is left closed and $\Psi$ right closed.
\end{itemize}
\end{itemize}

\subsection*{Section 2}

\begin{itemize}
\item $W \ \subset\ \Fl(\mathcal{M})$

\textit{$W$ weakly saturated} :
\begin{itemize}
\item[a)] $W$ contains the isomorphisms ;
\item[b)] if two among the three arrows of a commutative triangle are in $W$, the third too.
\end{itemize}

\item $W \ \subset\ \Fl(\mathcal{M})$

A morphism $p : X \longrightarrow Y$ of $\mathcal{M}$ is a \textit{$W$-fibration} ($p \in \mathrm{Fib}_W$) if for every diagram of cartesian squares
\end{itemize}